\chapter{Conclusion and Future Work}
\label{ch:conclusion}

\section{Summary of Contributions}

This thesis began from a puzzle: formal methods are mature and effective, yet used
only in a few niche corners of software engineering. Its diagnosis is that the
binding constraint is not the verification technology but the manual cost of the
specifications that technology requires, and its programme was to ask whether
automatic inference --- and the capabilities inferred specifications enable ---
could move formal methods out of the niche. It pursued the question through four
research questions, the first two asking whether inference removes the cost barrier
and the second two establishing the mechanisms by which library-boundary
specifications enable verification and compatibility analysis.

The answers, established empirically on production Java code, are as follows.

\emph{On inference accuracy (RQ1).} Java Modeling Language preconditions,
postconditions, and loop invariants can be inferred from source by pattern matching
organised around the weakest-precondition and strongest-postcondition calculi, at a
precision of 94.2\% for preconditions and 87.6\% for postconditions against
human-judged intent, with 92.9\% of all inferred clauses discharged by an
SMT-backed verifier against the implementation. The inference is heuristic and not
sound by construction; soundness is recovered \emph{a posteriori} by discharging
every clause and flagging the failures. The accuracy is high enough that automatic
inference is a credible substitute for the manual authoring that constitutes the
adoption barrier.

\emph{On downstream value (RQ2).} Cheaply inferred specifications deliver real
value: supplied as context to a large-language-model test generator, they improve
its unit-test generation by a large and statistically robust margin --- 272\% more
tests and a mutation score 40.7 percentage points higher than a signature-only
baseline. A guidance-only control rules out the alternative explanation that any
prompt enrichment would suffice, establishing that the specification is the active
ingredient. The cost reduction is therefore worth having: the cheaply produced
specifications are not cheap in value.

\emph{On library-boundary availability (RQ3).} Inferred JML contracts can be
carried with compiled Java libraries and service-interface descriptions, recovered
losslessly by a consumer without the source, at a bytecode-size overhead below
eleven percent and a throughput compatible with routine workflows. Across four
corpora exceeding twenty thousand methods the embedder achieves a 100\%
lossless-roundtrip rate through a toolchain that requires no modification. This
makes a dependency's contract available where verification of client code against
it must happen --- at the boundary, without the dependency's source.

\emph{On compositional propagation (RQ4).} A caller's specification depends
substantially on its callees': a top-down compositional analysis, measured across
five libraries and 21{,}052 methods, recovers 42{,}370 cross-method obligations the
single-pass propagation cannot express, dominated by the guard-conditional and
dispatch-dependent shapes through which a library change reaches the code that
calls it. This dependence is the mechanism that makes library-boundary
specifications a basis for behavioural-compatibility analysis: when a library's
inferred contracts change between versions, the impact on its callers' contracts is
computable through exactly these obligations. The thesis establishes the mechanism
and positions the end-to-end version-comparison experiment as its foremost future
work.

The thesis statement asserted that the adoption of formal methods is constrained
chiefly by the manual cost of specifications, that the constraint is substantially
removable by inference validated against a formal oracle, and that library-boundary
specifications --- because caller specifications derive from callee specifications
--- make verification across dependencies and behavioural-compatibility analysis
available to ordinary software. Each part is now backed by quantitative evidence on
real-world code.

\section{The Empirical Findings at a Glance}

It is worth gathering the principal quantitative findings in one place, because
their conjunction is the thesis's evidentiary core and they are otherwise
distributed across four chapters. On inference accuracy: the engine produces
non-trivial specifications for 99.0\% of the evaluation corpus, at 94.2\%
precondition precision and 87.6\% postcondition precision against documented intent,
with 92.9\% of all inferred clauses discharged by the verifier against the
implementation. On downstream value: specification-guided test generation produces
272\% more tests and a mutation score 40.7 percentage points higher than
signature-only generation, an effect size beyond conventional thresholds and not
reproduced by a guidance-only control. On library-boundary availability: across
four corpora exceeding twenty thousand methods, the embedder achieves a 100\%
lossless-roundtrip rate at a bytecode overhead below eleven percent and a
throughput of thousands of methods per second. On compositional propagation: a
caller's specification depends on its callees' through 42{,}370 cross-method
obligations recovered across five libraries and 21{,}052 methods, refining a
quarter to nearly a half of each library's methods, with the guard-conditional and
dispatch-dependent shapes --- those through which a library change reaches its
callers --- dominating.

These numbers are not independent assertions but a connected argument. The accuracy
figures establish that inference removes the cost barrier; the value figures
establish that the cheaply produced specifications are worth having; the
availability figures establish that they can be delivered to the library boundary
where verification of client code happens; and the propagation figures establish
that a library's specifications reach its callers' contracts strongly enough to
make behavioural-compatibility analysis a real prospect. Read together they answer
the adoption question the thesis posed: the specification-cost barrier is
removable, and removing it at the library boundary opens a capability ---
compatibility analysis across dependencies --- that the manual-specification regime
never made routine.

\section{Reflections on the Methodology}

Two methodological choices proved central enough to warrant restating as
lessons. The first is the two-layer architecture: infer aggressively by
heuristic pattern matching, validate conservatively by formal discharge. This
division --- a creative-but-fallible proposer and a rigorous-but-uncreative
checker --- is what allows the engine to be useful without being sound by
construction, and it is, in retrospect, the single design decision that makes
the rest of the thesis possible. An engine restricted to what it could prove
would emit too little; an engine with no formal back-stop would emit clauses no
consumer could trust. The separation is the resolution.

The second is the commitment to measurement on real code at scale. Every claim
in the thesis that could be made by counting --- clauses discharged, tests
generated, bytes overhead, clauses added by a second pass --- was made by
counting, on production libraries rather than on constructed examples. The
commitment is not merely a matter of rigour; it shaped what the thesis could
discover. A study built on constructed examples discovers what its author thought
to construct, and an author who believed per-spec compression would help, or that a
high roundtrip rate was good enough, or that more tests meant better tests, would
have constructed examples confirming those beliefs. Measurement on real artefacts
refused those beliefs: the compression made things worse, the roundtrip gap hid a
specific bug, the extra tests killed fewer mutants. The negative and surprising
results that the synthesis chapter collects are the dividend of the commitment, and
they are, in aggregate, as much a part of the thesis's contribution as the
confirmations, because they map the boundaries of the approach as precisely as the
confirmations map its successes. The cost
is that the conclusions are corpus-bound and quantitative rather than universal;
the benefit is that they are defensible, and that the failures they surface ---
the three tractability limits, the interface-method embedding bug, the
compile failures in the downstream probe --- are real failures of real systems
rather than artefacts of toy inputs. The thesis's most useful negative results,
including the counter-productive per-spec compression of
Chapter~\ref{ch:distribution} and the precise characterisation of what the
solver cannot discharge in Chapter~\ref{ch:testgen}, came directly from this
commitment.

\section{Limitations}

The principal limitations have been stated in the chapters and synthesised in
Chapter~\ref{ch:discussion}; they are collected here for completeness. The
evaluation corpora are uniformly mature, well-structured, open-source library
code, and the result that holds across them is correspondingly less established
on application code, framework code, and proprietary systems, and on the
event-handler and callback methods absent from the controlled corpus. The
language-model experiments use a single model family and defend an ordering of
conditions rather than absolute magnitudes. The formal validation is bounded by
the solver's tractability, leaving a residual class of clauses neither proved nor
refuted. And the downstream value of the compositional refinement rests, as yet,
on a test-count signal rather than on the mutation-kill metric the rest of the
thesis relies on. None of these undermines the central claims, but each bounds
their scope, and the honest summary is that the thesis establishes its
conclusions firmly on well-structured library code with a single model family,
and points beyond that scope rather than reaching it.

The limitations are also, read constructively, a map of where the most valuable
next work lies, and they are not uniform in their cost to address. The
single-model limitation is the cheapest to lift, requiring only the budget to re-run
parameterised experiments, and lifting it would most directly strengthen the
external validity of the value claim. The corpus-homogeneity limitation is more
expensive, requiring new corpora of application and framework code and, for some
method categories, new inference patterns, but it is the limitation whose lifting
would most broaden the thesis's reach. The residual-unverified-clauses limitation is
the hardest, depending on advances in SMT reasoning and JML translation that the
thesis does not itself deliver, and it bounds not the thesis's claims --- which
concern the verified clauses --- but the fraction of inferred clauses that can be
moved into the verified class. The compositional-downstream limitation is the most
immediately addressable of the empirical gaps, requiring only the repair of a test
harness and a mutation run on existing artefacts. Ordering the limitations by the
cost of addressing them is, in effect, the research roadmap the next section sets
out, and the fact that the cheapest-to-address limitations are also among the most
strengthening is the reason the thesis regards its conclusions as a foundation to
build on rather than a ceiling reached.

\section{Future Work}

The thesis opens several directions, in rough order of immediacy.

\paragraph{The version-to-version compatibility experiment.} The foremost
direction, because it would convert the thesis's central capability claim from a
demonstrated mechanism into a demonstrated result, is the version-comparison
experiment that Chapter~\ref{ch:compositional} positions as future work. The thesis
established that a caller's specification depends substantially on its callees', and
argued that this dependence makes a library change's impact on its callers'
contracts computable; the experiment that would close the argument runs the
inferrer over successive versions of a real library, recomputes the specifications
of client code against each version, and reports the behavioural incompatibilities
the comparison surfaces --- strengthened callee preconditions the client may now
violate, weakened callee postconditions the client relied on. Validated against
known breaking changes from the library's release history, such an experiment would
show that specification-level comparison detects behavioural incompatibilities that
signature compatibility misses, demonstrating the capability that is the thesis's
most direct answer to the adoption question. The mechanism is measured; the
end-to-end demonstration is the work that remains.

\paragraph{Completing the compositional downstream evaluation.} A second immediate
gap is the one named at the end of Chapter~\ref{ch:compositional}: the downstream
value of the more thoroughly propagated specifications is established only by a
test-count signal, because the generated tests did not all compile cleanly.
Completing the mutation-coverage comparison --- repairing the test-generation
harness so the generated suites compile and run, then measuring whether the
propagated specifications yield a higher mutation score --- would convert the
preliminary signal into a result on the same metric the rest of the thesis uses,
and would close the loop between Chapter~\ref{ch:compositional} and the value metric
of Chapter~\ref{ch:testgen}.

\paragraph{Service-boundary code.} The test-generation result of
Chapter~\ref{ch:testgen} was obtained on pure library code, where the oracle
problem is real but not at its most acute. The setting where it is most acute is
service-boundary code: methods that call REST endpoints or RPC stubs, where the
implementation delegates to an external service, the signature conveys nothing
about status codes or retry semantics, and signature-only generation collapses
into mocking the call. The distribution mechanism of Chapter~\ref{ch:distribution}
already extends to service interfaces via the OpenAPI extension, so the machinery
is in place to ask whether inferred specifications retain or amplify the
test-generation gain in that setting --- the question is whether the effect
reported here compounds where the oracle problem bites hardest.

\paragraph{Multi-model replication.} The single-model limitation is mechanical
to address and would strengthen every experiment that involves a language model.
The replication packages are structured for it --- the model client is
parameterised and the prompts are model-independent --- so the work is a matter
of budget rather than of design. A multi-model study would establish whether the
magnitudes of the effects, and not merely their direction, transfer across model
families.

\paragraph{Stronger loop-invariant inference.} The loop-invariant heuristics of
Chapter~\ref{ch:inference} produce a non-trivial invariant for roughly
eighty-five percent of loops, and the downstream consequence is measurable: the
methods left with only a weak invariant attain materially lower test quality.
Strengthening invariant inference --- through machine-learning approaches, or
through hybrid static-dynamic techniques in the spirit of dynamic invariant
detection --- would raise the floor on exactly the methods where the current
engine is weakest, and would attack the loop-related cases among the three
tractability limits.

\paragraph{Closing the verifier's tractability gaps.} The three tractability
limits of Chapter~\ref{ch:testgen} --- recursive multiplicative growth,
quantified invariants over mutable arrays, the for-each translation --- are
properties of the solver and the JML translation, not of the engine. Two of them
are plausibly addressable: the for-each translation gap is a question of exposing
the loop-variable-to-element relationship to user JML, and the mutable-array gap
is a question of array-theory formulation that further work on the OpenJML fork
might narrow. Closing them would move clauses from the unsupported class into the
verified class, shrinking the residual risk that the synthesis chapter
identified.

\paragraph{Single-pass capture of the easy compositional half.}
Chapter~\ref{ch:compositional} established that the branch-conditional half of
the compositional gap needs only local control-flow context, which a single-pass
analyser already has during its AST walk. An inferrer that wished to remain
single-pass could in principle lift its substituted clauses through enclosing
guards before adding them, capturing the dominant shape of the gap without the
machinery of a second pass. Whether this hybrid is worth building --- whether the
engineering simplicity of a single pass is better preserved by extending it or by
adding a clean second pass --- is a design question the measurement makes precise
but does not settle.

\paragraph{Beyond Java.} The system and its evaluation are Java-specific. The
WP/SP design discipline, the infer-then-verify architecture, and the
distribution mechanism are not intrinsically Java-bound, and porting the approach
to another language with a contract notation and a verifier --- or to a language
with neither, building both --- would test how much of the thesis's findings are
about specification inference in general and how much about Java in particular.

\paragraph{Hybrid static-dynamic and ensembled inference.} The thesis's engine is
purely static, and its complement --- dynamic invariant detection --- finds the
numeric and relational invariants the static engine cannot see. A hybrid that ran
both and took the union of their outputs, with the verifier filtering the combined
set, would cover more of the specification space than either alone, and the
infer-then-verify architecture extends naturally to it: both producers propose,
the verifier disposes. The same architecture admits a learned producer ---
prompting a model to propose specifications and filtering its proposals through the
verifier, in the spirit of Houdini with a model as the proposer --- which would
test whether the model's recall, combined with the verifier's soundness filter,
exceeds the static engine's precision-favouring output. The producer-consumer
asymmetry the thesis documents suggests where each producer is strong, and an
ensemble that played to each producer's strength is a direction the architecture
invites.

\section{A Research Roadmap}

The future directions above can be ordered into a roadmap by their dependencies and
their payoff. The highest-payoff step, requiring a new study but no new system, is
the version-to-version compatibility experiment: run the inferrer over successive
releases of a real library, recompute client specifications against each, and report
the behavioural incompatibilities the comparison surfaces, validated against the
library's known breaking changes. This is the step that would turn the thesis's
central capability claim --- that library-boundary specifications support
behavioural-compatibility analysis --- from a measured mechanism into a demonstrated
result, and it targets the adoption question most directly, because compatibility
analysis across dependencies is the capability that distinguishes inferred
library-boundary specifications from a merely cheaper version of what experts
already write.

A second, more immediate step, requiring no new technique, is to complete the
compositional downstream evaluation: repair the test-generation harness so the
generated suites compile, and measure whether the more thoroughly propagated
specifications raise the mutation score, converting the preliminary test-count
signal into a result on the thesis's primary metric. This closes the one empirical
gap the thesis leaves within its own scope.

A third step, requiring modest engineering, is the multi-model replication of
the test-generation experiments. The replication packages are structured for it,
and it would establish whether the magnitudes of the effects transfer across model
families or whether only their direction does --- a consequential strengthening of
the thesis's external validity, achievable without a new technique.

A fourth step, requiring a new study but no new system, is the service-boundary
evaluation. The distribution mechanism already reaches service interfaces through
the OpenAPI extension, so the machinery is in place; what is needed is a corpus of
service-boundary methods and a controlled experiment analogous to the library-code
one, testing the hypothesis that the specification's value compounds where the
source cannot serve as an oracle. This targets the setting where the oracle problem
is most acute.

A fifth step, requiring genuine new technique, is to attack the boundaries of
heuristic inference and of solver discharge --- the algorithmic postconditions the
static engine cannot reach, the recursive and array-quantified clauses the solver
cannot prove. Progress here is the hardest and least certain, depending on advances
in invariant inference and in SMT reasoning that the thesis does not itself
deliver, but it is the direction that would raise the ceiling on what the whole
approach can certify. The roadmap thus proceeds from demonstrating the thesis's
central capability, through completing its empirical gaps and strengthening its
external validity, to extending its reach and finally raising its ceiling --- a
sequence of increasing difficulty whose first step carries the highest payoff
because it makes the compatibility capability concrete.

\section{Broader Implications}

The findings of this thesis bear on three communities beyond the immediate one of
specification inference, and naming the implications situates the work in the
larger conversation it contributes to.

For the \emph{verification} community, the thesis offers a partial answer to the
specification-supply problem that has constrained the field's reach. Deductive
verifiers are powerful and mature, but they verify what they are given, and what
they are given is written by hand. An inference engine that supplies accurate
specifications automatically, even incomplete ones, changes the entry cost of
verification: a project can begin with machine-inferred contracts and refine the
ones that matter, rather than facing a blank slate. The discharge results of
Chapter~\ref{ch:testgen} suggest the inferred contracts are good enough that a
large majority pass the verifier unaltered, which means the human effort can
concentrate on the residual the machine cannot handle --- the recursive,
the quantified-over-arrays, the genuinely hard. The thesis does not claim to
automate verification; it claims to automate the specification step that
verification has always assumed, and to do so well enough that the assumption
becomes realistic.

For the \emph{testing} community, the thesis reframes the relationship between
specifications and test generation. The established tools solve the stimulus
problem and leave the oracle problem to the user; the thesis shows that the oracle
can be supplied automatically, from inferred specifications, and that doing so
improves a language-model generator's output by a large margin. The implication is
that the oracle problem and the specification-supply problem are the same problem
viewed from two directions, and that solving the second --- inferring
specifications --- solves the first --- supplying oracles --- as a consequence.
The finding that test count is a poor proxy for oracle quality, established by the
guidance control, is a methodological caution the testing community can apply
independently of the inference machinery: a test-generation technique should be
evaluated by mutation score or an equivalent fault-detection metric, not by the
volume of tests it produces.

For the community working at the intersection of \emph{language models and
software engineering}, the thesis offers a concrete instance of a general
principle: that learned models and grounded analyses are complementary, and that
the productive architecture uses each where it is strong. The model is a fluent
consumer of grounded information and an unreliable producer of it; the static
analysis is a reliable producer and no consumer at all. The systems that will get
the most from language models in software engineering, the thesis suggests, are
those that ground the model's fluency in analysis the model cannot perform ---
feeding it inferred specifications, verified facts, structured context --- rather
than those that ask the model to manufacture the grounding itself. The
producer-consumer asymmetry the thesis documents is a small, measured instance of
a pattern that recurs wherever correctness matters.

\section{Closing Remarks}

The barrier to formal methods, the thesis argued at the outset, is not the
verifier but the blank page --- the manual cost of writing the specifications
that the mature checking technology requires as input. This thesis has shown that
the blank page can be filled automatically, with specifications accurate enough
to be useful and compact enough to be distributed, and that the resulting
specifications carry measurable value to the contemporary tools, both classical
and learning-based, that consume them. It has also shown that the single-pass
inference that makes this practical leaves recoverable information behind, and
that a second pass recovers it.

The larger point is that inference, validation, distribution, and refinement
compose. A specification inferred but unverified cannot be trusted; a
specification verified but undistributed reaches no one; a specification
distributed but inferred only shallowly carries less than it could. The
contribution of this thesis is the demonstration that these stages can be made to
work together on real code at real scale --- that the path from unspecified Java
to verified, distributable, refined contracts can be walked automatically, and
that walking it is worth the effort because the contracts at the end of it
measurably improve the tools that read them. The manual cost that has kept formal
specifications on the margins of practice is, the evidence here suggests, no
longer the obstacle it was.

If the thesis has a single message beyond its individual results, it is that the
question worth asking about formal specifications has changed. For decades the
question was whether their benefit could justify their cost, and the answer, for
most software, was no --- the cost was too high and was paid by humans who had
other work. The thesis's evidence suggests the question is now different: not
whether the benefit justifies the cost, but what becomes possible when the cost is
largely removed. The studies here answer a first instalment of that question ---
better tests, reachable contracts, recoverable compositional information --- but
the question is larger than the answers, because a world in which every method has
an inferred, verified, distributed contract is a world in which a great deal of
tooling that has always assumed specifications as a scarce, manual input can assume
them as an abundant, automatic one. What that tooling does with abundant
specifications is a question the thesis opens rather than closes, and it is the
question the author hopes the work that follows will take up.
